\subsection{异步推理收益的实验验证}\label{sec:async-exp}

\subsubsection{实验目的与论证逻辑}\label{sec:async-exp-purpose}

4.4.5 节的不等式分析给出一个具体可证伪的预测：在当前
$T_{\text{infer}} \approx 2$ s、$N = 100$、$F = 30$ 的条件下，
异步推理理论稳态 FPS $\approx$ 20，与同步模式（约 19 FPS）几乎无差异。
本节设计实验对这一预测进行直接验证。

需要说明的是：实验的目的\textbf{不是证明异步推理普遍有效}，
而是\textbf{用实测数据为``异步推理在当前参数下不带来收益''
这一否定性结论提供工程证据}。这种\textbf{可行性边界的定量刻画}
本身具有工程价值——它能避免在错误的优化路径上投入工时，
也为``推理本体加速到什么程度后再启用异步''这一后续工作提供清晰的判据。

\subsubsection{实验设计}\label{sec:async-exp-design}

实验采用单变量对照，自变量为``调度模式'',
其余参数均严格控制以排除干扰：

\begin{center}
\begin{tabular}{lll}
  \toprule
  维度 & 取值 & 备注 \\
  \midrule
  调度模式 & \{同步, 异步\} & \textbf{唯一自变量} \\
  chunk\_size $N$ & 100 & 与第三章同步 sweep 对齐 \\
  $T_{\text{infer}}$ & $\approx 2$ s & 原始 ACT，不做量化或加速 \\
  目标 FPS $F$ & 30 & \texttt{environment\_dt = 33.3 ms} \\
  chunk\_size\_threshold & 0.5 & 异步默认值 \\
  任务 & 同一段录制好的回放任务 & 避免在线遥操污染 \\
  每模式重复 & $\geq$ 3 episode & 每段 $\geq$ 30 s，弃前 5 s warm-up \\
  \bottomrule
\end{tabular}
\end{center}

\subsubsection{控制变量}\label{sec:async-exp-control}

为保证两种模式之间的差异只来自调度方式，实验严格固定以下条件：
\begin{itemize}[leftmargin=2em]
  \item CPU 频率锁定为 \texttt{performance} governor，
        防止 DVFS 节流污染时序；
  \item \texttt{torch.set\_num\_threads(3)} 与 \texttt{cv2.setNumThreads(1)}
        在两种模式下保持一致；
  \item 异步模式下 client 和 server 同机部署，
        通过 \texttt{taskset} 把 server 绑核到 0--2、client 绑核到 3，
        与 4.4.6 节描述一致；
  \item 每个 episode 前后用 \texttt{vcgencmd measure\_temp} 记录 CPU 温度，
        记录每秒平均频率，用于排除热降频；
  \item Git commit、模型权重、相机分辨率与采集格式
        （MJPEG 1280$\times$720 @30 FPS）保持一致。
\end{itemize}

\subsubsection{测量指标}\label{sec:async-exp-metrics}

为充分验证 4.4 节理论分析，本实验在 client 与 server 两端
插入逐帧时间戳打点，采集六类指标：

\begin{center}
\begin{tabular}{p{3.2cm}p{5.2cm}p{4cm}}
  \toprule
  类别 & 指标 & 采集方法 \\
  \midrule
  吞吐 & \texttt{fps\_mean}, \texttt{fps\_p5}, \texttt{fps\_p95}
       & 沿用 \texttt{timing\_stats.csv} 字段 \\
  推理耗时 & $T_{\text{infer}}$ 实测分布（mean, std）
       & \texttt{policy\_server} 端 \texttt{predict\_action\_chunk()} 前后打点 \\
  流水线健康 & stall 帧数比例、stall 平均时长、\texttt{must\_go} 触发次数
       & \texttt{robot\_client} 主循环逐帧记录 \\
  Chunk 利用率 & 每 chunk 过期帧 / 实际执行帧比例
       & 用 \texttt{latest\_action} 与 chunk \texttt{timestep} 计算 \\
  系统资源 & 4 核 CPU 占用率、CPU 温度漂移、CPU 平均频率
       & \texttt{psutil} + \texttt{vcgencmd} 每秒采样 \\
  端到端延迟 & 观测采集 $\to$ 对应 action 执行的时间差
       & \texttt{TimedObservation} / \texttt{TimedAction} 时间戳差 \\
  \bottomrule
\end{tabular}
\end{center}

其中 \textbf{stall 帧比例}与 \textbf{chunk 过期帧比例}是支撑 4.4 节
不等式分析的关键实证：若理论分析正确，异步模式下应同时观察到显著的
stall 比例（情况 C）和约 60\% 的过期帧比例（$T_{\text{infer}} F / N = 60/100$）。

\subsubsection{实验流程}\label{sec:async-exp-procedure}

\textbf{步骤 1：环境准备}。固定 git commit 与模型权重 hash，
锁定 CPU governor，关闭 WiFi 等可能引入抖动的后台服务，
散热风扇全速。

\textbf{步骤 2：模型预热}。两种模式各先跑 5 个 episode 不计入数据，
让 PyTorch 算子缓存、相机自动曝光、舵机温度补偿全部进入稳态。

\textbf{步骤 3：同步基线采集}。直接运行 \texttt{record.py}，
沿用第三章 \texttt{timing\_stats.csv} 的采集脚本，得到
\texttt{obs\_s}、\texttt{inference\_s}、\texttt{action\_s}、\texttt{wait\_s}
四阶段耗时。

\textbf{步骤 4：异步实验采集}。两个终端分别启动：
\begin{verbatim}
# 终端 1（policy_server，绑核 0-2）
taskset -c 0-2 python -m lerobot.scripts.server.policy_server \
  --host 127.0.0.1 --port 50051 ...
# 终端 2（robot_client，绑核 3）
taskset -c 3 python -m lerobot.scripts.server.robot_client \
  --server_address 127.0.0.1:50051 \
  --chunk_size_threshold 0.5 ...
\end{verbatim}
两端各加日志钩子，按帧打点
（\texttt{timestep, latest\_action, queue.qsize, must\_go, T\_infer}）。

\textbf{步骤 5：数据对齐}。实验后用 \mytodo{}
脚本按 \texttt{timestep} 把 client 和 server 日志做 join，
得到每帧的 \texttt{obs\_send\_t}、\texttt{action\_recv\_t}、\texttt{action\_exec\_t}，
进而计算端到端延迟与 chunk 过期帧比例。

\textbf{步骤 6：统计分析}。每模式 $\geq$ 3 episode 取 mean $\pm$ std，
使用配对 t 检验（按 episode 配对）评估同步与异步的差异是否显著。

\subsubsection{预期结果与待验证假设}\label{sec:async-exp-hypotheses}

基于 4.4 节的理论分析，本实验预期得到以下四条结果。
论文实验数据填入后将逐条核对：

\begin{itemize}[leftmargin=2em]
  \item \textbf{H1（吞吐）}：
        $|$ \texttt{fps\_mean}(异步) $-$ \texttt{fps\_mean}(同步) $|$ $\leq$ 10\%。
        即异步在当前参数下\textbf{不优于}同步，
        吻合不等式预测的稳态 FPS $\approx$ 20。
  \item \textbf{H2（stall）}：异步模式下 stall 比例显著大于 0，
        且呈现以推理周期为单位的周期性出现，
        与图~\ref{fig:async-inf-three-cases} 的情况 C 拓扑一致。
  \item \textbf{H3（chunk 利用率）}：异步模式下每 chunk 约 60\% 帧数被丢弃为过期帧，
        与理论值 $T_{\text{infer}} F / N = 60/100$ 定量吻合。
  \item \textbf{H4（CPU 争抢）}：异步模式下树莓派 5 实测 $T_{\text{infer}}$
        相对 server 独占 4 核时抬高约 10--25\%（绑核策略可缓解但不能完全消除）。
\end{itemize}

H2 和 H3 是\textbf{不等式分析的直接定量证据}。
若它们与预测显著不符（例如 stall 比例 $<$ 5\%、或过期帧比例 $<$ 30\%），
则应回到 4.4.4 节重新审视推导前提；
若完全符合预测，则可作为``异步推理在当前条件下尚不可用''的实证锚点。

\subsubsection{实验结果呈现（待回填）}\label{sec:async-exp-results}

实验完成后，本节将以以下五项结果为核心呈现：
\begin{enumerate}[leftmargin=2em]
  \item \textbf{图：FPS 对比柱状图}（同步 vs 异步，含误差线）；
  \item \textbf{图：稳态时间线甘特图}，
        实测画出执行 / 推理两条 lane，stall 段高亮染色，
        与图~\ref{fig:async-inf-three-cases} 情况 C 作对比；
  \item \textbf{图：chunk 利用率堆叠柱图}，
        展示每 chunk 中``已过期被丢弃 / 实际执行''的帧数分解；
  \item \textbf{表：8 项关键指标主结果表}（fps、stall\%、过期\%、CPU\%、
        端到端延迟、$T_{\text{infer}}$、\texttt{must\_go} 触发频次、CPU 温度漂移）；
  \item \textbf{表：配对 t 检验显著性}（同步 vs 异步在 fps 上的 $p$ 值）。
\end{enumerate}

\mytodo{实验未跑，结果数据与图表待回填；以上图表骨架为最低呈现要求。}

\subsubsection{讨论：实验结论的工程含义}\label{sec:async-exp-discussion}

预期实验将印证 4.4.5 节的理论预测：在当前 $T_{\text{infer}} = 2$ s 下，
异步推理的工程收益接近于零。然而这并不意味着异步推理本身没价值——
它意味着\textbf{异步推理的可用性强依赖于推理本体的速度}。
当 $T_{\text{infer}}$ 通过量化、ONNX 或 NEON 优化下降到
\begin{equation}
  T_{\text{infer}} \;\leq\; \frac{N}{2F} \;=\; \frac{100}{2 \times 30} \;\approx\; 1.67\ \text{s}
\end{equation}
以下时，不等式重新成立，异步推理才会从``无效''切换为``有效''。
也就是说，在本平台上，\textbf{异步推理是``推理加速完成后''的下游红利，
而不是独立可拿的优化}。这一判据将直接支撑 4.6 节的优化路线图。
